\section{Background}
\subsection{IO500}

\subsubsection{Motivation and purpose}

High Performance Computing (HPC) systems are used to run applications
that process large amounts of data. As a result, the performance of the
storage system and I/O operations can have a significant impact on the
overall performance of an application.

However, evaluating I/O performance is not straightforward. Different
storage systems can have different architectures and configurations, and
their performance can vary depending on the workload and the way the
system is used. A common and reproducible methodology is therefore
needed to compare the I/O performance of different systems.

The IO500 benchmark suite was created to address this need. It provides
a common methodology for evaluating the I/O performance of HPC storage
systems using several workloads and a defined set of execution rules.
The benchmark is developed together with the HPC community and is based
on existing open-source benchmarks~\cite{io500about}.

The results obtained with IO500 can be used to compare different storage
systems and configurations. They are also collected in public IO500
lists, which provide a common reference for the HPC community. This
makes IO500 useful not only for measuring performance, but also for
studying how different systems and configurations affect I/O performance.

In this work, IO500 is used as the target benchmark for studying
automated configuration tuning. The objective is to determine whether a
search strategy can find a high-performing configuration while reducing
the number of benchmark evaluations required.

\subsubsection{Benchmark structure}

IO500 is not a single benchmark program. It is a benchmark suite that
combines several workloads under a common set of execution and
submission rules. The suite uses existing open-source benchmarks and
defines how they should be executed so that results can be compared
across systems~\cite{io500about}.

The benchmark covers different aspects of storage performance. In
particular, it includes workloads targeting data transfer performance,
metadata operations, and file-system search operations. Each workload
produces an individual performance measurement. These measurements are
then combined to produce an overall IO500 score.

The IO500 website publishes the results of submitted runs together with
information about the tested system. This information includes, among
other elements, the storage system, file system, number of client nodes,
number of processes, and measured performance. The results are organised
into several lists, including Production, Research, and Full lists
\cite{io500list}.

The benchmark therefore provides both individual measurements and an
overall score. Keeping the individual results is useful because two
systems with similar overall scores may have very different performance
characteristics.

The following subsection describes the main workloads used by IO500 and
the type of I/O operation measured by each of them.
\subsection{IOPS}
\subsection{Bayesian Optimization}